% Includable fragment. The construction was proposed by the root research
% agent and independently checked here. Requires the F/D/epsilon setting
% and theorem environments from the parent paper.
\section{Coverage and a right predecessor without successor injectivity}
\label{sec:right-predecessor}

There is a stronger consequence of the self-seeded equations than finite
aperiodicity. After reinstating contextual induction on literal equations,
the source can prove numerical coverage and construct a finite-case
right predecessor, still without successor injectivity. The word
``right'' is deliberate: the result is $S(P(Sn))=Sn$, not $P(Sn)=n$.

Use the computation, equality-test, and separation axioms stated in
Section~\ref{sec:finite-aperiodicity}, together with the contextual
literal-equation induction rule. Write $D(n,a)=F_{v.a}(n)$ as before.
All functions introduced below, except the explicitly graph-defined $P$,
are finite source terms.

\subsection{Quantifier-free induction is available without cancellation}

For Boolean values define
\[
 \neg_B b:=\epsilon(b,0),\qquad
 b\land_B c:=\epsilon\bigl(D(b,\epsilon(c,0)),0\bigr),\qquad
 b\lor_B c:=\neg_B(\neg_B b\land_B\neg_B c).
\]
Their truth tables use only $D(0,a)=Sa$, $D(1,a)=a$, equality reflection,
and $0\ne1$. In the conjunction table, the sole non-Boolean intermediate
value is $S1$; its distinction from $0$ follows from zero separation,
without requiring $S1\ne1$ or cancelling a successor.

External recursion on a fixed quantifier-free formula $A$ therefore
constructs a Boolean term $\chi_A$ with $\chi_A=1\leftrightarrow A$.
Implication is represented by $\neg_B\chi_A\lor_B\chi_B$; this is valid
intuitionistically because the finite equality characteristics make the
quantifier-free antecedent decidable. Literal-equation induction on
$\chi_A=1$ consequently yields contextual quantifier-free induction.
No successor-injectivity instance is used in this conversion.

\subsection{A clock that stops one step before its target}

For a passive parameter $a$, define
\[
 b_a(v):=\epsilon(v,0)\lor_B\epsilon(Sv,a),\qquad
 h_a(v):=D(b_a(v),v),\qquad
 V(j,a):=F_{h_a}(j),\qquad W(a):=V(a,a).
\]
The Boolean guard gives the exact two cases
\[
 h_a(v)=v\quad\text{if }v=0\text{ or }Sv=a,\qquad
 h_a(v)=Sv\quad\text{otherwise}.
\]
In particular $h_a(0)=0$, so the self-seeded base is
\[
 V(0,a)=1,\qquad V(Sj,a)=h_a(V(j,a)).
\]

\begin{lemma}[Capped-clock invariant]
\label{lem:rp-clock}
The source proves, for all $j,a$,
\[
                       V(j,a)=Sj\quad\lor\quad S(V(j,a))=a.
\]
\end{lemma}
\begin{proof}
Use quantifier-free induction on $j$, keeping $a$ fixed. The base is
$V(0,a)=1=S0$. For the step put $v=V(j,a)$. If the second alternative
$Sv=a$ holds, the update fixes $v$, preserving that alternative. Otherwise
the old first alternative gives $v=Sj$, so $v\ne0$. The guard is then
false, and the update gives $Sv=S(Sj)$, the first alternative at $Sj$.
Equivalently, one may split the Boolean guard first; a successful guard
in the first alternative cannot be caused by $v=0$ and therefore gives
$Sv=a$. Every case uses only the displayed induction hypothesis and
zero/successor separation.
\end{proof}

\begin{lemma}[Right predecessor on double successors]
\label{lem:rp-double}
For every $m$, the source proves
\[
                         S(W(S^2m))=S^2m.
\]
\end{lemma}
\begin{proof}
First use Lemma~\ref{lem:rp-clock} with its passive parameter still
independent, and only then specialize $a=S^2m$ and $j=m$. In its first
alternative $V(m,a)=Sm$; its second alternative already states
$S(V(m,a))=a$. Thus both alternatives prove the latter equation.
The guard therefore fixes $V(m,a)$. Two explicit successor computation
steps give
\[
 V(S^2m,a)=h_a(h_a(V(m,a)))=V(m,a).
\]
Since $a=S^2m$, the left side is $W(S^2m)$. Applying $S$ establishes
the conclusion. No successor is cancelled.
\end{proof}

\subsection{A finite graph and a legitimate induction instance}

Define $P$ by the finite quantifier-free graph
\[
\begin{split}
 \Phi_P(n,p):={}&(n=0\land p=0)\ \lor\ (n=1\land p=0)\\
 &{}\lor\ (n\ne0\land n\ne1\land p=W(n)).
\end{split}
\]
Equality decidability and $0\ne1$ prove this graph total and single
valued. Introduce $P$ in its conservative definitional extension.
This use of $P$ does not enlarge the primitive iteration schema to new
updates containing $P$.

For the induction needed next, finite-case elimination supplies more than
mere conservativity. Each term with finitely many $P$ occurrences has a
finite partition into source quantifier-free guards with source value
terms: refine cases for its arguments, and then use the three displayed
cases of $\Phi_P$. Repeating this finite expansion translates any
quantifier-free formula of the extension into an equivalent source
quantifier-free formula. Consequently contextual quantifier-free induction
is admissible in this extension. No general principle that definitional
extensions preserve restricted induction is being assumed.

\begin{theorem}[Coverage and a right predecessor are derived]
\label{thm:rp-coverage}
The source, through the finite graph above, proves
\[
 P0=P1=0,\qquad
 n=0\ \lor\ n=S(Pn),\qquad
 S(P(Sn))=Sn.
\]
It consequently proves $n=0\lor\exists p\;n=Sp$, without a coverage
axiom or successor-injectivity axiom.
\end{theorem}
\begin{proof}
The two boundary equations are immediate from the graph. Use the
legitimate quantifier-free induction just described on
$n=0\lor n=S(Pn)$.

At zero the first alternative holds. For the successor step, if the old
first alternative $n=0$ holds, then $Sn=1=S(P1)$. In the old second
alternative, $Sn=S^2(Pn)$, so Lemma~\ref{lem:rp-double} gives
$S(W(Sn))=Sn$. The input $Sn$ is nonzero. If $Sn=1$, its graph value
is zero and $S(P(Sn))=1=Sn$. If $Sn\ne1$, its graph value is $W(Sn)$,
and the same conclusion follows from the lemma. This proves the
successor instance of the induction formula in every case.

Apply the resulting coverage formula to $Sn$ and exclude its zero
alternative to obtain $S(P(Sn))=Sn$. Existential introduction gives the
ordinary coverage statement. All uses of $P$ can be removed by the
finite graph translation.
\end{proof}

\paragraph{The remaining cancellation problem.}
The theorem chooses one predecessor of each nonzero element. It does not
prove that every predecessor of that element is the chosen one. Equating
$Sa$ and $Sb$ permits applying $P$, but only gives $P(Sa)=P(Sb)$; the
missing statements are $P(Sa)=a$ and $P(Sb)=b$. Thus this construction
removes coverage from the no-injectivity precursor theory, but does not
yet remove the injectivity axiom from the existing complete iPRA
presentation.

\begin{lemma}[An exact form of the remaining closure problem]
\label{lem:rp-retraction}
Put $T(n):=P(Sn)$ and $C(n):\!\iff T(n)=n$ in the finite-case
definitional extension. Then
\[
 S(T(n))=Sn,\qquad T(T(n))=T(n),\qquad T(0)=0.
\]
Over the precursor theory, successor injectivity is equivalent to
\[
                         \forall n\;(C(n)\to C(Sn)).
\]
\end{lemma}
\begin{proof}
The first equation is Theorem~\ref{thm:rp-coverage}. Applying $P$ to it
gives the second; the third is $P1=0$. Thus $T$ is a retraction onto
its fixed points, with $0$ among them.

If these fixed points are successor closed, the already justified
quantifier-free induction on $C(n)$ gives $T(n)=n$ for every $n$.
From $Sa=Sb$, application of $P$ then gives $a=b$. Conversely,
successor injectivity applied to $S(T(n))=Sn$ gives $T(n)=n$ for
every $n$, hence the displayed closure condition.
\end{proof}

\paragraph{Why a direct comparison of the clocks is insufficient.}
One might try to establish this closure by comparing the updates with
targets $Sn$ and $S^2n$. For $v\ne0$, the expression $S(h_{Sn}(v))$
uses its stopping branch when $Sv=Sn$, whereas $h_{S^2n}(Sv)$ uses
its stopping branch when $S^2v=S^2n$. The putative intertwining
identity therefore needs
\[
                         S^2v=S^2n\ \Longrightarrow\ Sv=Sn
\]
at least for every state visited by the clocks. In the contrary case
the two sides are $S^2v$ and $Sv$, which are distinct by fixed-point
exclusion. The hypothesis $C(n)$ specifies the chosen predecessor of
$Sn$; no proved invariant turns it into uniqueness at $S^2n$.
This identifies a gap in that attempted proof, not an independence
result.
